\subsection{Energy Model for 5G NR gNB and UEs}

To evaluate the network energy consumption during simulations, the ns-3 energy framework is extended by implementing custom \texttt{NrGnbEnergyModel} and \texttt{NrUeEnergyModel} modules \cite{samorzewski_energy}.
These models track the operational states of the NR spectrum PHY layer, represented by
the state space $\mathcal{S} = \{\mathrm{IDLE}, \mathrm{RX\_CTRL}, \mathrm{RX\_DATA}, \mathrm{TX}, \mathrm{CCA\_BUSY}\}$.
The power draw of each node is computed from the time spent in each state.

For the User Equipment (UE), the total energy consumed is computed from its base power profile:
\begin{equation}
    E_\mathrm{UE} = \int_{0}^{T} P_\mathrm{UE}\!\left(s(t)\right)\, dt = \sum_{s \in \mathcal{S}} P_s^\mathrm{UE} \cdot \tau_s,
\end{equation}
where $P_s^\mathrm{UE}$ is the power consumption in state $s$, and $\tau_s$ is the cumulative time spent in that state.

By contrast, the Next Generation NodeB (gNB) power model accounts for both a baseline operating load and the scaling effects of its antenna array:
\begin{equation}
    E_\mathrm{gNB} = \int_{0}^{T} P_\mathrm{gNB}(s(t))\, dt = \sum_{s \in \mathcal{S}} (P_\mathrm{fixed} + N_\mathrm{ant} P_s^\mathrm{gNB}) \tau_s,
\end{equation}
where $P_\mathrm{fixed}$ is the constant hardware power draw, $N_\mathrm{ant}$ is the number of antennas, and $P_s^\mathrm{gNB}$ is the per-antenna state power.

The time a node spends in active states ($\tau_s$) depends on the network's link adaptation, which relies on the Signal-to-Interference-Plus-Noise Ratio (SINR).
For a receiver $i$ on subcarrier $k$, the SINR is calculated as:
\begin{equation}
    \mathrm{SINR}_{i,k} = \frac{P_\mathrm{tx}\, G_\mathrm{tx}\, G_\mathrm{rx}\, |H_{i,k}|^2}{N_0 W_k + \sum_{j \neq i} P_\mathrm{tx}\, G_\mathrm{tx}\, G_\mathrm{rx}\, |H_{j,k}|^2},
\end{equation}
where $H$ is the channel coefficient supplied by the propagation model described in Section~\ref{subsec:propagation-calculation}.

The gNB uses this SINR to select a Modulation and Coding Scheme (MCS); a higher SINR enables a more spectrally efficient MCS, which reduces the time spent in the power-intensive transmit state ($\tau_{\mathrm{TX},i}$).
Overall energy efficiency, measured as energy per delivered bit, is then:
\begin{equation}
    \mathcal{E}_b = \frac{E_\mathrm{node}}{\sum_i \eta_i\, W_i\, \tau_{\mathrm{TX},i}},
\end{equation}
where $E_\mathrm{node}$ is the total energy of the UE or gNB.

This formulation shows that energy efficiency $\mathcal{E}_b$ is highly sensitive to the channel coefficient $H$.
A simple analytical model such as Friis provides a coarse approximation, whereas ray tracing yields an environment-specific, geometry-accurate channel.
Experiment~II evaluates how much this choice of propagation model alters the final energy metrics.
